\documentclass[11pt,a4paper]{article}

\usepackage[utf8]{inputenc}
\usepackage[T1]{fontenc}
\usepackage{lmodern}
\usepackage[french]{babel}
\usepackage[margin=2.4cm]{geometry}
\usepackage{microtype}
\usepackage{booktabs}
\usepackage{array}
\usepackage[table]{xcolor}
\usepackage{amsmath}
\usepackage{amssymb}
\usepackage{mathtools}
\usepackage{enumitem}
\usepackage[most]{tcolorbox}
\usepackage{caption}
\usepackage{titlesec}
\usepackage{hyperref}

\definecolor{navy}{HTML}{1F3864}
\definecolor{bluemid}{HTML}{2E5496}
\definecolor{lightblue}{HTML}{D6E0F0}
\definecolor{accent}{HTML}{C0491F}

\hypersetup{colorlinks=true, linkcolor=bluemid, urlcolor=bluemid,
  pdftitle={Draft - Probabilites conditionnelles au sein d'un comptage N_ij},
  pdfauthor={Kelian Kaddouri}}

\newtcolorbox{cle}[1][]{colback=lightblue!40,colframe=navy,boxrule=0.6pt,
  arc=2pt,left=8pt,right=8pt,top=6pt,bottom=6pt,before skip=9pt,after skip=9pt,#1}
\newtcolorbox{attention}[1][]{colback=accent!8,colframe=accent,boxrule=0.6pt,
  arc=2pt,left=8pt,right=8pt,top=6pt,bottom=6pt,before skip=9pt,after skip=9pt,#1}

\titleformat{\section}{\normalfont\large\bfseries\color{navy}}{\thesection}{0.6em}{}
\titlespacing*{\section}{0pt}{1.9ex plus .3ex minus .2ex}{1.0ex plus .2ex}
\captionsetup{font=small,labelfont={bf,color=navy},labelsep=period,skip=6pt}

\title{\vspace{-1.2cm}\color{navy}\bfseries Draft : probabilités conditionnelles au sein d'un comptage $N_{i,j}$\\[2pt]
  \large De la cascade qualitative à la brique fréquence/sévérité d'un SCR}
\author{Kélian Kaddouri}
\date{\today}

\begin{document}
\maketitle
\vspace{-0.5cm}

\begin{cle}
\textbf{But de ce draft (document de travail).} On note $N_{i,j}$ le \textbf{nombre d'incidents}
observés dans la cellule (segment~$i$, pilier~$j$) : c'est une \textbf{fréquence}. L'objet
d'étude est la loi \emph{conditionnelle} à ce comptage : \og{}sachant qu'il y a eu $N_{i,j}$
incidents, comment chacun se propage-t-il à travers les piliers~?\fg{} La réponse est que la
\textbf{cascade qualitative devient la loi conditionnelle de propagation}, et cette brique
alimentera ensuite la sévérité, la perte agrégée, puis le SCR. Le SCR lui-même n'est pas traité
ici : c'est l'étape d'après.
\end{cle}

\section{Le cadre : fréquence, puis sévérité}

On décompose la perte d'une cellule en deux étages, comme tout \textbf{modèle collectif de
risque}~\cite{Buhlmann1970,KlugmanPanjerWillmot,Mikosch2009}.
\begin{itemize}[leftmargin=1.4em,itemsep=1pt]
\item \textbf{Fréquence} : $N_{i,j}$, le nombre d'incidents entrant par le pilier~$j$ dans le
  segment~$i$ sur la période. C'est un comptage aléatoire. Pour le cyber, on le prendra
  \emph{surdispersé} (un \textbf{Poisson mixte}~\cite{Grandell1997} plutôt qu'un Poisson simple :
  un choc commun fait varier l'intensité, ce qui est le pendant fréquentiel du facteur
  systémique et rejoint les modèles de contagion cyber~\cite{BessyRolandBoumezouedHillairet2021,HillairetLopez2021}).
\item \textbf{Sévérité conditionnelle} : \emph{sachant} $N_{i,j}=n$, chacun des $n$ incidents
  produit un dommage. C'est \textbf{ici} que vit \og{}au sein d'un même $N_{i,j}$\fg{} : la loi
  des dommages \emph{conditionnelle au comptage}.
\end{itemize}

\begin{cle}
\textbf{Le lien avec la cascade.} Un incident n'est pas un point : il \emph{se propage}. Entré
par le pilier~$j$, il peut entraîner d'autres piliers. La \textbf{sévérité d'un incident est
donc gouvernée par sa cascade}, et la cascade se décrit par des \textbf{probabilités
conditionnelles} $\mathbb P(j' \mid j)$. C'est exactement l'objet que le modèle qualitatif
fournit déjà, une fois traduit en probabilités.
\end{cle}

\section{La probabilité conditionnelle de propagation $\mathbb P(j' \mid j)$}

On part de la matrice dirigée TRANS ($\text{TRANS}[j][j']$ = force \og{}$j$ entraîne $j'$\fg{}).
Ce n'est pas encore une probabilité : les lignes ne somment pas à~1. On \textbf{normalise chaque
ligne} par sa somme $s_j = \sum_{j'} \text{TRANS}[j][j']$ :
\begin{equation}
\boxed{\;\mathbb P(j' \mid j) \;=\; \frac{\text{TRANS}[j][j']}{s_j}\;,\qquad
\sum_{j'\neq j} \mathbb P(j'\mid j) = 1.\;}
\end{equation}
C'est la probabilité que \textbf{le prochain pilier touché soit $j'$, sachant que $j$ vient de
l'être}. Les $\mathbb P(j'\mid j)$ forment une \textbf{matrice de transition} : la cascade
devient une \textbf{chaîne de Markov} sur les piliers~\cite{Norris1997}, dont la structure
dirigée reprend les modèles de propagation d'influence~\cite{Pearl1988,Kempe2003}. L'asymétrie
causale est \emph{préservée}.

\begin{table}[h]\centering\small
\renewcommand{\arraystretch}{1.3}
\begin{tabular}{@{}c|ccccc|c@{}}
\toprule
$\mathbb P(j'\mid j)$ & $j'=$P1 & P2 & P3 & P4 & P5 & $s_j$ \\
\midrule
$j=$P1 & ---    & $0{,}308$ & $0{,}269$ & $0{,}269$ & $0{,}154$ & $2{,}60$ \\
P2     & $0{,}133$ & ---    & $0{,}267$ & $0{,}200$ & $0{,}400$ & $1{,}50$ \\
P3     & $0{,}214$ & $0{,}357$ & ---    & $0{,}214$ & $0{,}214$ & $1{,}40$ \\
P4     & $0{,}118$ & $0{,}471$ & $0{,}176$ & ---    & $0{,}235$ & $1{,}70$ \\
P5     & $0{,}143$ & $0{,}286$ & $0{,}286$ & $0{,}286$ & ---    & $0{,}70$ \\
\bottomrule
\end{tabular}
\caption{Le noyau de propagation, chaque ligne normalisée à~1. Exemple d'asymétrie :
$\mathbb P(\text{P2}\mid\text{P1})=0{,}31$ mais $\mathbb P(\text{P1}\mid\text{P2})=0{,}13$. La
colonne $s_j$ retient la force d'émission de chaque pilier.}
\end{table}

\section{L'état d'arrêt : indispensable pour une sévérité finie}

\begin{attention}
\textbf{Sans arrêt, la cascade ne s'éteint jamais.} Si chaque pilier touché en entraîne toujours
un autre (lignes sommant à~1), l'incident se propage sans fin et sa sévérité diverge. Il faut un
\textbf{état d'arrêt} $\varnothing$ (un état \emph{absorbant}) pour que le nombre de piliers
atteints soit fini, exactement la condition d'extinction d'un processus de
branchement~\cite{AthreyaNey1972}.
\end{attention}

\noindent On introduit une \textbf{propension à propager} $e_j = g\cdot s_j / \max_k s_k$, avec un
gain $g\in(0,1)$ (ici $g=0{,}9$, $\max_k s_k = s_1 = 2{,}60$), et l'on pose
\begin{equation}
\mathbb P(j' \mid j) = e_j \cdot \frac{\text{TRANS}[j][j']}{s_j},
\qquad
\mathbb P(\varnothing \mid j) = 1 - e_j.
\end{equation}
Un pilier qui émet peu éteint souvent la cascade : on garde \emph{où} \textbf{et} \emph{si} elle
continue. Ce gain $g$ est le même objet que la normalisation de la partie quantitative du
mémoire (il borne la propagation, comme un $R_0<1$).

\begin{table}[h]\centering\small
\renewcommand{\arraystretch}{1.25}
\begin{tabular}{@{}cccccc@{}}
\toprule
pilier $j$ & P1 & P2 & P3 & P4 & P5 \\
\midrule
propension $e_j = 0{,}9\,s_j/2{,}60$ & $0{,}90$ & $0{,}52$ & $0{,}48$ & $0{,}59$ & $0{,}24$ \\
$\mathbb P(\text{arrêt}\mid j) = 1-e_j$ & $0{,}10$ & $0{,}48$ & $0{,}52$ & $0{,}41$ & $\mathbf{0{,}76}$ \\
\bottomrule
\end{tabular}
\caption{Avec l'état d'arrêt, P5 éteint la cascade dans $76\,\%$ des cas, P1 rarement : la force
d'émission est restaurée, et toute cascade se termine presque sûrement.}
\end{table}

\section{Ce que \og{}au sein d'un même $N_{i,j}$\fg{} donne concrètement}

Sachant $N_{i,j}=n$ incidents dans la cellule, \textbf{chacun} des $n$ incidents part du
pilier~$j$ et réalise une cascade tirée du noyau ci-dessus. La probabilité d'un chemin est un
produit de probabilités conditionnelles. Exemple, incident entré en P1 :
\begin{align}
\mathbb P(\text{P1 seul}) &= \mathbb P(\varnothing\mid\text{P1}) = 0{,}10, \\
\mathbb P(\text{P1}\to\text{P2}\to\varnothing) &= \underbrace{e_1\tfrac{0{,}80}{2{,}60}}_{\mathbb P(\text{P2}\mid\text{P1})=0{,}277}\times \underbrace{\mathbb P(\varnothing\mid\text{P2})}_{0{,}48} = 0{,}133 .
\end{align}
La \textbf{sévérité} de l'incident se déduit des piliers atteints, via GBASE
(P1$=6$, P2$=4$, P3$=3$, P4$=7$, P5$=1$) : \og{}P1 seul\fg{} coûte $6$ ; \og{}P1$\to$P2\fg{}
coûte $6+4=10$. La loi conditionnelle des sévérités est donc \emph{entièrement déterminée} par
les probabilités conditionnelles de propagation.

\begin{cle}
\textbf{La perte de la cellule, et le pont vers le SCR.} La perte agrégée est une
\textbf{somme composée}~\cite{Buhlmann1970,KlugmanPanjerWillmot} : sachant $N_{i,j}=n$, on somme
les sévérités des $n$ cascades,
\[
S_{i,j} \;=\; \sum_{k=1}^{N_{i,j}} X_k, \qquad
X_k = \text{sévérité de la cascade de l'incident } k .
\]
Les $X_k$ ont la loi conditionnelle décrite ici. \textbf{L'étape suivante} (pas ce draft) est
la loi de $S_{i,j}$, puis l'agrégation sur les cellules et la VaR/TVaR à $99{,}5\,\%$ : le
SCR~\cite{mcneil2015,SolvencyII2009}.
\end{cle}

\section{Sachant quels piliers ont fait défaut : la probabilité de l'ordre}

\begin{cle}
\textbf{La question.} On \emph{sait} que $m$ piliers donnés ont fait défaut, par exemple
$S=\{\text{P1},\text{P2},\text{P3}\}$. On ne sait pas dans quel \textbf{ordre} la défaillance
s'est propagée. On cherche la \textbf{probabilité conditionnelle de chaque ordre, sachant
l'ensemble} $S$. Il y a $m!$ ordres possibles ($3!=6$ pour trois piliers) : c'est exactement le
comptage $N_m = m\cdot N_{m-1}$.
\end{cle}

À chaque ordre on associe un \textbf{poids}, égal à la probabilité de ce chemin dans le modèle :
la propension à \emph{amorcer} sur le premier pilier (ROOT), multipliée par les forces
d'enchaînement dirigées (TRANS) le long de la chaîne,
\begin{equation}
W(j_1 \to j_2 \to \dots \to j_m) \;=\; \text{ROOT}[j_1]\,\prod_{\ell=1}^{m-1}\text{TRANS}[j_\ell][j_{\ell+1}].
\end{equation}
La probabilité conditionnelle d'un ordre est son poids \textbf{divisé par la somme des poids de
tous les ordres} du même ensemble (le conditionnement \og{}sachant que ce sont ces piliers\fg{}
se lit dans le dénominateur, qui ne parcourt que les permutations de $S$) :
\begin{equation}
\boxed{\;\mathbb P\big(j_1 \to \dots \to j_m \;\big|\; \text{défaut de } S\big)
\;=\; \frac{W(j_1 \to \dots \to j_m)}{\displaystyle\sum_{\sigma \in \mathfrak S(S)} W(\sigma)}\;.}
\end{equation}

\noindent\textbf{Exemple chiffré, $S=\{\text{P1},\text{P2},\text{P3}\}$.} Avec
$\text{ROOT}=\{1{,}0\,;0{,}6\,;0{,}5\}$ pour (P1, P2, P3) et la matrice TRANS du modèle :

\begin{table}[h]\centering\small
\renewcommand{\arraystretch}{1.3}
\begin{tabular}{@{}lccc@{}}
\toprule
Ordre & Poids $W$ (ROOT $\times$ TRANS $\times$ TRANS) & Poids & Probabilité \\
\midrule
\rowcolor{lightblue!40} P1 $\to$ P3 $\to$ P2 & $1{,}0 \times 0{,}70 \times 0{,}50$ & $0{,}350$ & $\mathbf{35{,}1\,\%}$ \\
P1 $\to$ P2 $\to$ P3 & $1{,}0 \times 0{,}80 \times 0{,}40$ & $0{,}320$ & $32{,}1\,\%$ \\
P3 $\to$ P1 $\to$ P2 & $0{,}5 \times 0{,}30 \times 0{,}80$ & $0{,}120$ & $12{,}0\,\%$ \\
P2 $\to$ P1 $\to$ P3 & $0{,}6 \times 0{,}20 \times 0{,}70$ & $0{,}084$ & $8{,}4\,\%$ \\
P2 $\to$ P3 $\to$ P1 & $0{,}6 \times 0{,}40 \times 0{,}30$ & $0{,}072$ & $7{,}2\,\%$ \\
P3 $\to$ P2 $\to$ P1 & $0{,}5 \times 0{,}50 \times 0{,}20$ & $0{,}050$ & $5{,}0\,\%$ \\
\midrule
\multicolumn{2}{r}{somme des poids} & $0{,}996$ & $100\,\%$ \\
\bottomrule
\end{tabular}
\caption{Probabilité conditionnelle de chaque ordre, sachant que P1, P2, P3 ont fait défaut.
On calcule les six poids, puis on divise par leur somme pour ramener le total à $100\,\%$.}
\end{table}

\begin{cle}
\textbf{Lecture.} Sachant que ces trois piliers ont fait défaut, l'ordre le plus probable est
\textbf{P1 $\to$ P3 $\to$ P2} ($35{,}1\,\%$). Les deux ordres qui démarrent par P1 totalisent
$67\,\%$ : c'est cohérent, P1 est le plus fort amorceur (ROOT $=1{,}0$) et le plus fort émetteur.
On répond ainsi précisément à la question \og{}dans quel ordre~?\fg{} : le modèle fournit la
\textbf{loi de probabilité sur les ordres}, pas un ordre unique.
\end{cle}

\begin{attention}
\textbf{Précision.} Ce calcul suppose que l'on conditionne sur \emph{ces piliers exactement}, et
compare leurs $m!$ ordres entre eux : les probabilités d'arrêt et de non-propagation vers les
autres piliers, communes à la comparaison, disparaissent au rapport. Si l'on veut la probabilité
\emph{jointe} \og{}exactement ces piliers, dans cet ordre, et aucun autre\fg{}, il faut réintégrer
l'état d'arrêt de la section précédente ; la loi sur les ordres, elle, reste celle du tableau.
\end{attention}

\section{Le jour où l'on aura des données}

Les scores TRANS jouent aujourd'hui le rôle des comptages. Le jour où un registre fournit les
transitions observées, la probabilité conditionnelle s'estime \emph{de la même façon}, par le
comptage normalisé, l'estimateur classique des transitions d'une chaîne de
Markov~\cite{AndersonGoodman1957} :
\[
\widehat{\mathbb P}(j' \mid j) \;=\; \frac{N_{j,j'}}{\sum_{k} N_{j,k}} .
\]
La structure du modèle ne change pas ; seules les entrées deviennent empiriques.

\begin{attention}
\textbf{Questions ouvertes à trancher.}
\begin{enumerate}[leftmargin=1.5em,itemsep=1pt,topsep=2pt]
\item Loi de fréquence pour $N_{i,j}$ : Poisson, ou mixte Poisson / binomiale négative pour la
  surdispersion cyber~?
\item Garde-t-on l'état d'arrêt (sévérité finie) et le gain $g$, ou une autre borne de
  propagation~?
\item La sévérité par pilier : GBASE brut, ou la formule de gravité complète du modèle
  qualitatif (étendue $+$ amplification)~?
\item Dépendance entre cellules $N_{i,j}$ d'un même segment $i$ (le choc commun) : à modéliser
  pour l'agrégation SCR.
\end{enumerate}
\end{attention}

\section{Sources des notions mobilisées}

Le draft assemble des briques actuarielles et probabilistes standard. Chaque notion et sa source :

\begin{table}[h]\centering\small
\renewcommand{\arraystretch}{1.3}
\begin{tabular}{@{}p{6.6cm} p{6.6cm}@{}}
\toprule
\textbf{Notion} & \textbf{Source} \\
\midrule
Modèle collectif fréquence/sévérité & Bühlmann~(1970)~\cite{Buhlmann1970} ; Klugman, Panjer \& Willmot~\cite{KlugmanPanjerWillmot} ; Mikosch~(2009)~\cite{Mikosch2009} \\
Somme composée $S=\sum_{k=1}^{N} X_k$ & Bühlmann~(1970)~\cite{Buhlmann1970} ; Klugman et coll.~\cite{KlugmanPanjerWillmot} \\
Fréquence Poisson mixte (surdispersion) & Grandell~(1997)~\cite{Grandell1997} \\
Choc commun / contagion fréquentielle cyber & Bessy-Roland et coll.~(2021)~\cite{BessyRolandBoumezouedHillairet2021} ; Hillairet \& Lopez~(2021)~\cite{HillairetLopez2021} \\
Probabilité conditionnelle, chaîne de Markov & Norris~(1997)~\cite{Norris1997} \\
Propagation dirigée (TRANS $\to$ transition) & Pearl~(1988)~\cite{Pearl1988} ; Kempe et coll.~(2003)~\cite{Kempe2003} \\
État absorbant / extinction de cascade & Athreya \& Ney~(1972)~\cite{AthreyaNey1972} \\
Estimation $\widehat{\mathbb P}(j'|j)=N_{j,j'}/N_{j,\cdot}$ & Anderson \& Goodman~(1957)~\cite{AndersonGoodman1957} \\
VaR/TVaR à $99{,}5\,\%$, SCR, agrégation & McNeil, Frey \& Embrechts~(2015)~\cite{mcneil2015} ; Directive Solvabilité~II~\cite{SolvencyII2009} \\
\bottomrule
\end{tabular}
\end{table}

\begin{thebibliography}{99}

\bibitem{Buhlmann1970}
Bühlmann, H.
\emph{Mathematical Methods in Risk Theory}. Springer, 1970.

\bibitem{KlugmanPanjerWillmot}
Klugman, S.~A., Panjer, H.~H. et Willmot, G.~E.
\emph{Loss Models: From Data to Decisions}. Wiley (éditions successives).

\bibitem{Mikosch2009}
Mikosch, T.
\emph{Non-Life Insurance Mathematics: An Introduction with the Poisson Process}.
2\ieme\ éd., Springer, 2009.

\bibitem{Grandell1997}
Grandell, J.
\emph{Mixed Poisson Processes}. Chapman \& Hall, 1997.

\bibitem{BessyRolandBoumezouedHillairet2021}
Bessy-Roland, Y., Boumezoued, A. et Hillairet, C.
\og Multivariate Hawkes process for cyber insurance\fg{}.
\emph{Annals of Actuarial Science}, 15(1):14--39, 2021.

\bibitem{HillairetLopez2021}
Hillairet, C. et Lopez, O.
\og Propagation of cyber incidents in an insurance portfolio\fg{}.
\emph{Scandinavian Actuarial Journal}, 2021.

\bibitem{Norris1997}
Norris, J.~R.
\emph{Markov Chains}. Cambridge University Press, 1997.

\bibitem{Pearl1988}
Pearl, J.
\emph{Probabilistic Reasoning in Intelligent Systems: Networks of Plausible Inference}.
Morgan Kaufmann, 1988.

\bibitem{Kempe2003}
Kempe, D., Kleinberg, J. et Tardos, {\'E}.
\og Maximizing the Spread of Influence through a Social Network\fg{}.
\emph{Proc. 9th ACM SIGKDD}, p.~137--146, 2003.

\bibitem{AthreyaNey1972}
Athreya, K.~B. et Ney, P.~E.
\emph{Branching Processes}. Springer, 1972.

\bibitem{AndersonGoodman1957}
Anderson, T.~W. et Goodman, L.~A.
\og Statistical Inference about Markov Chains\fg{}.
\emph{The Annals of Mathematical Statistics}, 28(1):89--110, 1957.

\bibitem{mcneil2015}
McNeil, A.~J., Frey, R. et Embrechts, P.
\emph{Quantitative Risk Management: Concepts, Techniques and Tools}.
éd. révisée, Princeton University Press, 2015.

\bibitem{SolvencyII2009}
Parlement européen et Conseil de l'Union européenne.
\emph{Directive 2009/138/CE (Solvabilité~II)}. 2009.

\end{thebibliography}

\end{document}
